% Part: normal-modal-logic
% Chapter: axioms-systems
% Section: normal-logics

\documentclass[../../../include/open-logic-section]{subfiles}

\begin{document}

\olfileid[pt]{nml}{prf}{nor}

\olsection{Lógicas modais normais}

Nem todo conjunto de !!{formula}s modais pode ser facilmente caracterizado
como o conjunto das !!{formula}s deriváveis de um conjunto de axiomas.
Queremos que as lógicas modais se comportem bem. Em primeiro lugar, tudo o que
podemos !!{derive} na lógica proposicional clássica deve continuar sendo
!!{derivable}, levando naturalmente em conta que as !!{formula}s agora também
podem conter \iftag{prvBox}{$\Box$\iftag{prvDiamond}{
    e~}{}}{}\iftag{prvDiamond}{$\Diamond$}{}. Para isso, exigimos que uma
lógica modal contenha todas as instâncias tautológicas e seja fechada sob
modus ponens.

\begin{defn}
  Uma \emph{lógica modal} é um conjunto~$\Sigma$ de !!{formula}s modais que
  \begin{enumerate}
  \item contém todas as tautologias, e
  \item é fechado sob substituição, isto é, se $!A \in \Sigma$, e
    $!D_1$, \dots, $!D_n$ são !!{formula}s, então
    \[
    \SSubst{!A}{\subst{!D_1}{p_1}, \dots, \subst{!D_n}{p_n}} \in \Sigma,
    \]
    \item é fechado sob \emph{modus ponens}, isto é, se $!A$ e
      $!A \lif !B \in \Sigma$, então $!B \in \Sigma$.
  \end{enumerate}
\end{defn}

Para usar a semântica relacional das lógicas modais, também precisamos exigir
que todas as !!{formula}s válidas em todos os modelos modais estejam
incluídas. Ocorre que essa exigência é satisfeita assim que todas as instâncias
de \Ax{K}\iftag{prvDiamond}{ e de~\Dual{}}{} são !!{derivable}, e sempre que
!!a{formula}~$!A$ é !!{derivable}, também o é~$\Box !A$. Uma lógica modal que
satisfaz essas condições é chamada de \emph{normal}. (Naturalmente, também
existem lógicas modais não normais, mas os modelos relacionais usuais não são
adequados para elas.)

\begin{defn}
  Uma lógica modal $\Sigma$ é \emph{normal} se contém
  \begin{align*}
    \tag{\Ax{K}} & \Box(p \lif q) \lif (\Box p \lif \Box q),
    \iftag{prvDiamond}{\\
      \tag{\Dual} & \Diamond p \liff \lnot\Box\lnot p}{}
  \end{align*}
  e é fechada sob \emph{necessitação}, isto é, se $!A \in \Sigma$, então
  $\Box !A \in \Sigma$.
\end{defn}

Observe que, enquanto a implicação tautológica tem uma granularidade
suficientemente fina para preservar a \emph{verdade em um mundo}, a regra
\Nec{} preserva apenas a \emph{verdade em um modelo} (e, portanto, também a
validade em um frame ou em uma classe de frames).

\begin{prop}\ollabel{prop:rk}
  Toda lógica modal normal é fechada sob a regra \RK,
  \begin{prooftree}
    \AxiomC{$!A_1 \lif (!A_2 \lif \cdots (!A_{n-1} \lif !A_n)\cdots)$}
    \RightLabel{\RK}
    \UnaryInfC{$\Box!A_1 \lif (\Box!A_2 \lif \cdots (\Box!A_{n-1}
      \lif \Box!A_n)\cdots).$}
  \end{prooftree}
\end{prop}

\begin{proof}
  Por indução sobre~$n$: se $n = 1$, então a regra é simplesmente \Nec, e
  toda lógica modal normal é fechada sob \Nec.

  Suponha agora que o resultado valha para $n-1$; mostramos que vale para~$n$.

  Suponha
  \begin{align*}
  & !A_1 \lif (!A_2 \lif \cdots (!A_{n-1} \lif !A_n)\cdots) \in \Sigma
  \intertext{Pela hipótese de indução, temos}
  & \Box!A_1 \lif (\Box!A_2 \lif \cdots \Box(!A_{n-1} \lif !A_n)\cdots)
  \in \Sigma
  \intertext{Como $\Sigma$ é uma lógica modal normal, ela contém todas as
    instâncias de~$\Ax{K}$, em particular}
  & \Box(!A_{n-1} \lif !A_n) \lif (\Box!A_{n-1} \lif \Box!A_n) \in \Sigma
  \intertext{Usando modus ponens e instâncias tautológicas adequadas, obtemos}
  & \Box!A_1 \lif (\Box!A_2 \lif \cdots (\Box!A_{n-1}
  \lif \Box!A_n)\cdots) \in \Sigma. 
  \end{align*}
\end{proof}

\begin{prop}\ollabel{prop:notDiamondBot}
  Toda lógica modal normal $\Sigma$ contém~$\lnot\Diamond\lfalse$.
\end{prop}

\begin{prob}
  Demonstre \olref[nml][prf][nor]{prop:notDiamondBot}.
\end{prob}

\begin{prop}
  Sejam $!A_1$, \dots, $!A_n$ !!{formula}s. Então existe uma menor lógica
  modal $\Sigma$ que contém todas as instâncias de $!A_1$, \dots, $!A_n$.
\end{prop}

\begin{proof}
  Dados $!A_1$, \dots, $!A_n$, defina $\Sigma$ como a interseção de todas as
  lógicas modais normais que contêm todas as instâncias de $!A_1$, \dots,
  $!A_n$. A interseção é não vazia, pois $\Frm[L]$, o conjunto de todas as
  !!{formula}s, é uma lógica modal desse tipo.
\end{proof}

\begin{defn}
A menor lógica modal normal que contém $!A_1$, \dots, $!A_n$ é chamada de
\emph{sistema modal} e denotada por $\Log{K} !A_1 \dots !A_n$. A menor lógica
modal normal é denotada por~\Log{K}.
\end{defn}

\end{document}
